\documentclass[12pt]{article}
\usepackage[margin=0.52in]{geometry}
\usepackage{graphicx}
\usepackage{booktabs}
\usepackage{microtype}
\usepackage{enumitem}
\setlength{\parindent}{0pt}
\setlength{\parskip}{3pt}
\setlist{nosep,leftmargin=*}
\newif\ifpendingmeasurements
\pendingmeasurementstrue
\newcommand{\pending}[1]{\textit{[Student: insert #1 from verified measurements]}}
\newcommand{\plotplaceholder}[2]{%
  \IfFileExists{#1}{\includegraphics[width=#2]{#1}}%
    {\fbox{\parbox[c][1.2in][c]{#2}{\centering PENDING MEASUREMENTS\\#1}}}}
\begin{document}
\ifpendingmeasurements
  \begin{center}\fbox{\large\textbf{PENDING MEASUREMENTS — NOT A SUBMISSION REPORT}}\end{center}
\fi
\begin{center}
{\Large CS 425 MP3: Distributed Key-Value Store}\par
Krishna \pending{NetID} and Juhi \pending{NetID} \quad University of Illinois Urbana-Champaign
\end{center}

\textbf{Design} (keep this section to at most one page.)
\pending{concise description of fixed five-owner SHA-256 placement, MP2 membership,
durable SQLite records, Strong two-phase ABD, Eventual R/W, Eventual-RR repair,
restart recovery, and the correctness evidence actually obtained}.

\textbf{LLM use} (fewer than 300 words.)
\pending{truthful summary of assistance used and human verification/design work.
Preserve complete project transcripts under \texttt{llm\_use/}; do not claim a transcript
exists until the team has saved it.}

\newpage
\textbf{Performance} (Experiment 1.) Strong and Eventual-noRR at R/W=1/1 and 2/2;
N=5, ten VMs, 1 KiB values, 100,000 preloaded keys, uniform key selection, 50\% GET
and 50\% PUT, with 1, 2, 4, 8 and 10 clients. Each point is the mean of three
independent trials of at least 180 seconds per client; error bars show sample standard
deviation across trial summaries.

\begin{center}\plotplaceholder{exp1_throughput_vs_latency.png}{0.88\linewidth}\end{center}
\pending{brief discussion of observed throughput/latency trends, variability, errors,
and whether results match the team's expectations, based only on the plotted trials}.

\textbf{Performance with read repair} (Experiment 2.) Eventual-noRR versus Eventual-RR
at R/W=1/1, with ten clients and 99\% of operations directed to the same ten hot keys.
The three independent trials per configuration lasted at least 180 seconds per client.

\begin{center}
\plotplaceholder{exp2_throughput.png}{0.47\linewidth}\hfill
\plotplaceholder{exp2_mean_latency.png}{0.47\linewidth}
\end{center}
\pending{brief measured comparison of throughput and latency, including variability,
hot-key effects, failures, and a cautious explanation of observed read-repair cost}.

\end{document}
